% TMLR submission skeleton.
% For actual TMLR submission: download tmlr.sty from
% https://github.com/JmlrOrg/tmlr-style-file, replace the documentclass block with
%   \documentclass{article}\usepackage[preprint]{tmlr}
% and adjust author formatting to the TMLR template. This file compiles standalone
% with the plain article class so the content can be iterated without the style file.

\documentclass[11pt]{article}
\usepackage[margin=1in]{geometry}
\usepackage{booktabs}
\usepackage{graphicx}
\usepackage{amsmath}
\usepackage{natbib}
\usepackage{hyperref}
\usepackage{xcolor}
\newcommand{\todo}[1]{\textcolor{red}{[TODO: #1]}}

\title{How Much of the Elliptic Leaderboard Is Real?\\
Decomposing Evaluation Inflation into Leakage, Shift, and Base-Rate Effects\\
on Temporal Transaction Graphs}

\author{Julien Audibert\\ \texttt{[affiliation]}}
\date{}

\begin{document}
\maketitle

\begin{abstract}
The Elliptic Bitcoin dataset and its heterogeneous extension Elliptic++ are standard
benchmarks for graph-based fraud detection, and recent work reports near-perfect
illicit-detection performance (F1 up to $\sim$0.98). We show these numbers are largely
an artifact of evaluation protocol. Using a single codebase, we train the same models
(GCN, GraphSAGE, GAT, XGBoost) under (a) the honest temporal split and (b) a random
split, and find the random split inflates illicit-class PR-AUC by $+0.195\pm0.001$
(XGBoost) to $+0.553\pm0.055$ (GCN) over five seeds; under the random split a vanilla
XGBoost reaches F1 0.955 / PR-AUC 0.986, squarely in the reported ``SOTA'' band. We
then \emph{decompose} the gap with two controls. A prevalence-matched random test set
isolates a small base-rate component ($+0.026$ GNN / $+0.005$ XGBoost). An inductive
ablation, which removes all test nodes from the graph during training, measures the
graph-specific ``message-passing leakage'' we ourselves initially hypothesized at
\textbf{zero} ($-0.000\pm0.007$). The dominant component is \emph{distribution
access} ($+0.384$ GNN / $+0.189$ XGBoost): random splits train models on the test
period's distribution, and the models that generalize worst across temporal shift
(GNNs) gain the most. The inflation generalizes to the Elliptic++ address task under
a deployment-honest feature protocol ($+0.497$), while a temporally stable
cross-domain control (DGraph-Fin) shows essentially no PR-AUC inflation ($+0.002$),
consistent with the shift-driven mechanism. Under the only comparable protocol,
gradient-boosted trees still lead graph neural networks, reproducing the original
2019 benchmark; a 12-configuration GNN tuning sweep narrows but does not close the
gap. We characterize the unsolved post-shift regime and report four novel but
negative drift-robust method attempts, including the first adaptation of USAD's
two-decoder adversarial objective to graphs. All results are reproducible from a
CI-gated public repository.
\end{abstract}

\section{Introduction}
Detecting illicit activity in blockchain transaction graphs is a canonical
application of graph machine learning. The Elliptic dataset
\citep{weber2019aml} (203{,}769 Bitcoin transactions across 49 time steps,
$\sim$2\% illicit) and Elliptic++ \citep{elmougy2023elliptic} (adding 822{,}942
wallet addresses) are the de facto benchmarks. A stream of recent papers reports
illicit-detection F1 approaching 0.98, suggesting the problem is nearly solved. We
argue this impression is false and traces to two evaluation choices: (i) random
rather than temporal train/test splits, which give models access to the test
period's distribution, and (ii) reporting weighted or overall F1, dominated by the
$\sim$98\% licit class, rather than illicit-class metrics.

Our contributions are: (1) a systematic, multi-model, multi-seed quantification of
the random-split inflation on Elliptic-family benchmarks; (2) a decomposition of
that inflation into base-rate, graph-leakage, and distribution-access components,
including the measured refutation of our own message-passing-leakage hypothesis;
(3) a temporally stable cross-domain control consistent with the shift-driven
mechanism; and (4) honest negative results for four drift-robust method variants,
including a first USAD-on-graph adaptation. We release a leakage-free evaluation
harness in which every reported number is one seeded command.

\section{Related work}
\paragraph{Evaluation integrity.} Leakage was formalized by \citet{kaufman2011leakage}
(the ``no time machine'' rule); \citet{kapoor2023leakage} taxonomize it as a driver
of ML reproducibility failures; \citet{arp2022dos} name ``temporal snooping'' among
security-ML pitfalls; \citet{shchur2018pitfalls} document protocol sensitivity in
GNN evaluation generally. The closest template is TESSERACT
\citep{pendlebury2019tesseract}, which showed Android-malware F1 of $\sim$0.99
collapsing under time-aware evaluation. Relative to this line, our contribution is
graph-native and mechanistic: we quantify the inflation on Elliptic-family
benchmarks across model classes and decompose it, showing the graph-specific
channel is null and the dominant term is distribution access.

\paragraph{Elliptic-specific evaluation.} \citet{weber2019aml} established the
temporal protocol and the trees-over-GNNs result; GADBench \citep{tang2023gadbench}
confirms tuned tree ensembles lead on Elliptic (tuned RF AUPRC $\approx$0.789,
matching our XGBoost). A recent preprint \citep{maganti2026liability} critiques
Elliptic evaluation along a transductive-versus-inductive axis. None of these
measure random-split inflation across models, decompose it, or test a
stable-domain control.

\paragraph{Graph anomaly detection and USAD.} Reconstruction autoencoders
\citep{ding2019dominant}, generative and adversarially regularized variants, and
one-class methods form the unsupervised graph-anomaly literature. USAD
\citep{audibert2020usad}, a two-decoder adversarially trained autoencoder for
multivariate time series, has to our knowledge not previously been adapted to
graphs; we adapt it and report the (negative) outcome.

\section{Experimental protocol}
\label{sec:protocol}
All experiments use one codebase, fixed seeds, and a CI-gated environment.
The honest protocol is: a temporal split (Elliptic's canonical early-versus-late
time steps; first-seen splits for addresses and DGraph-Fin), illicit-class PR-AUC
and F1 as primary metrics, a validation slice carved from the latest train steps,
and decision thresholds chosen on validation only. Leakage comparisons hold the
model, features, and trainer fixed and change only the split. Models: GCN
\citep{kipf2017gcn}, GraphSAGE \citep{hamilton2017graphsage}, GAT
\citep{velickovic2018gat}, XGBoost \citep{chen2016xgboost}.

\section{Results}
\subsection{The honest baseline ordering}
Under the temporal split, XGBoost (PR-AUC $0.791\pm0.001$; F1 0.82 under the
full-train protocol) leads GraphSAGE ($0.510\pm0.035$), GAT ($0.297\pm0.054$) and
GCN ($0.265\pm0.051$), reproducing \citet{weber2019aml}. A 12-configuration tuning
sweep selected on temporal validation improves GraphSAGE to PR-AUC 0.506 / F1 0.486
(from F1 0.422 at defaults) without changing the ordering. Our GNN F1 remains below
some honest-protocol literature values ($\sim$0.69); we report this openly as a
protocol/training-detail discrepancy rather than claiming parity.

\subsection{The inflation, quantified}
\begin{table}[h]
\centering
\caption{Illicit-class PR-AUC under the honest temporal split versus a stratified
random split (mean $\pm$ std over 5 seeds). Same model, same trainer, same
features; only the split differs.}
\label{tab:inflation}
\begin{tabular}{lccc}
\toprule
Model & Temporal & Random & Inflation \\
\midrule
GCN & $0.265\pm0.051$ & $0.818\pm0.009$ & $+0.553\pm0.055$ \\
GraphSAGE & $0.510\pm0.035$ & $0.920\pm0.007$ & $+0.410\pm0.041$ \\
GAT & $0.297\pm0.054$ & $0.775\pm0.032$ & $+0.478\pm0.036$ \\
XGBoost & $0.791\pm0.001$ & $0.986\pm0.001$ & $+0.195\pm0.001$ \\
\bottomrule
\end{tabular}
\end{table}

Under the random split, XGBoost reaches F1 0.955 / PR-AUC 0.986, inside the
literature's reported $\sim$0.90--0.98 band. The effect generalizes to the
Elliptic++ address task (heterogeneous graph, 822k address nodes) under a
deployment-honest feature protocol (pre-split wallet aggregation, train-only
scaling): over three seeds, PR-AUC $0.417\pm0.042$ (temporal) versus
$0.966\pm0.002$ (random), an inflation of $+0.548\pm0.041$.

\subsection{Decomposing the gap}
\label{sec:decompose}
Two controls, threshold-free PR-AUC, five seeds. (i) Matching the random test
set's prevalence to the temporal test's (6.5\%) isolates the base-rate component.
(ii) An inductive ablation removes all test nodes from the graph during training,
so no test feature or test-adjacent edge is visible; this measures graph-specific
message-passing leakage directly.

\begin{table}[h]
\centering
\caption{Decomposition of the temporal-versus-random gap (PR-AUC, 5 seeds).}
\label{tab:decompose}
\begin{tabular}{lcccc}
\toprule
Component & GCN & GraphSAGE & GAT & XGBoost \\
\midrule
Base rate & $+0.061\pm0.004$ & $+0.026\pm0.005$ & $+0.073\pm0.011$ & $+0.005\pm0.001$ \\
Message-passing leakage & $+0.003\pm0.016$ & $-0.000\pm0.007$ & $-0.010\pm0.033$ & n/a \\
Distribution access & $+0.489\pm0.042$ & $+0.384\pm0.044$ & $+0.415\pm0.044$ & $+0.189\pm0.002$ \\
\midrule
Total gap & $+0.553\pm0.055$ & $+0.410\pm0.041$ & $+0.478\pm0.036$ & $+0.195\pm0.001$ \\
\bottomrule
\end{tabular}
\end{table}

The message-passing channel we ourselves hypothesized is \textbf{null across all
three architectures} ($+0.003\pm0.016$ GCN, $-0.000\pm0.007$ GraphSAGE,
$-0.010\pm0.033$ GAT): training with test nodes entirely removed from the graph
leaves random-split PR-AUC unchanged. The dominant term everywhere is distribution
access: a random split trains the model on examples drawn from the test period, an
advantage no deployed model can have. GNNs inflate more than XGBoost because
their temporal generalization is weaker, not because of a graph-specific leak.
We report this as an explicit self-correction.

\subsection{A temporally stable control}
On DGraph-Fin \citep{huang2022dgraph} (3.7M nodes, timestamped edges, temporally
stable $\sim$1.3\% fraud rate), the same GraphSAGE shows essentially no PR-AUC
inflation ($0.037 \rightarrow 0.039$, $+0.002$), consistent with the mechanism:
where the distribution does not shift, there is nothing to access. A wider model
(hidden dimension 128) behaves identically (PR-AUC inflation $+0.004$; ROC-AUC
$+0.055$). Caveats stated plainly: ROC-AUC does inflate ($+0.031$ to $+0.055$),
both arms sit near the PR-AUC floor at this model strength, and GADBench reports
strong GNNs on DGraph-Fin, so this control is supporting evidence rather than
proof.

\subsection{The unsolved post-shift regime and negative method results}
Under the temporal protocol, a rolling per-window backtest of EvolveGCN-O
\citep{pareja2020evolvegcn} collapses at the dark-market shutdown (steps 44--46:
PR-AUC $\sim$0.01) and only partially recovers. A learning-rate $\times$
gradient-clipping sweep (six configurations) improves the aggregate test PR-AUC
from 0.100 to at best 0.200 (lr $5{\times}10^{-3}$, no clipping), still far below
every honest baseline; the conclusion is unchanged.
Four drift-robust variants fail honestly: a USAD-style two-decoder adversarial
graph autoencoder \citep{audibert2020usad} (naive, rolling-normal, and
domain-adversarial variants, all PR-AUC $\sim$0.037) and a supervised
domain-adversarial GraphSAGE \citep{ganin2015dann} (PR-AUC 0.364, below the plain
baseline 0.510): enforcing time-invariance removes signal rather than improving
generalization under this regime change.

\section{Discussion and recommendations}
(1) On temporally shifting fraud graphs, report illicit-class metrics under a
strict temporal split; random-split numbers can overstate PR-AUC by up to
$+0.55$. (2) The post-shift window is the real benchmark; aggregate scores hide
the collapse. (3) Graph structure is not a free win on Elliptic: features already
encode neighborhood aggregates, and under honest evaluation trees lead. (4)
Method claims should be accompanied by decomposition-style controls; our own
double-leak hypothesis did not survive its ablation.

\section{Conclusion}
The near-perfect reported performance on Elliptic-family benchmarks is largely an
evaluation artifact whose dominant mechanism is distribution access, not
graph-specific leakage. We provide the decomposition, the controls, a reproducible
harness, and honest negative results, and we re-anchor the honest state of the art.

\bibliographystyle{plainnat}
\bibliography{references}

\end{document}
